\begin{figure}[!htbp]
\centering
\begin{minipage}[t]{.31\linewidth}
\centering
\resizebox{\linewidth}{!}{%
\begin{tikzpicture}[scale=2.1,>=Latex]
  \node[panellabel,anchor=west] at (-1.08,1.18) {(a) two CE2 gaps};
  \HexCoords
  \DrawHexagon
  \DrawRays
  \draw[gap] ($(V0)!.18!(V1)$)--($(V0)!.36!(V1)$);
  \draw[gap] ($(V5)!.64!(V0)$)--($(V5)!.84!(V0)$);
  \draw[centerrole]
    (.07,.03)--(.45,.54)--(.62,-.35)--cycle;
  \node[figlabel,text=GapPurple] at (.55,.69) {demand $q$};
  \node[figlabel,text=GapPurple] at (.78,-.58) {demand $p$};
  \node[figlabel,align=center] at (0,-1.18)
    {$p,q$ are the common\\boundary demands};
\end{tikzpicture}%
}
\end{minipage}\hfill
\begin{minipage}[t]{.37\linewidth}
\centering
\resizebox{\linewidth}{!}{%
\begin{tikzpicture}[x=5cm,y=1cm,>=Latex]
  \node[panellabel,anchor=west] at (-.02,2.45) {(b) transverse center exits};
  \draw[flow] (0,1.55)--(1,1.55);
  \draw[flow] (0,.35)--(1,.35);
  \fill (0,1.55) circle (.018) node[left=2pt,figlabel] {$O$};
  \fill (1,1.55) circle (.018) node[right=2pt,figlabel] {$V_2$};
  \fill (0,.35) circle (.018) node[left=2pt,figlabel] {$O$};
  \fill (1,.35) circle (.018) node[right=2pt,figlabel] {$V_4$};
  \coordinate (E2) at (.08,1.55);
  \coordinate (D2) at (.14,1.55);
  \coordinate (D4) at (.14,.35);
  \coordinate (E4) at (.19,.35);
  \fill[CenterBlue] (E2) circle (.024) node[above=3pt,figlabel] {C trace end};
  \fill[DemandRed] (D2) circle (.024) node[below=3pt,figlabel] {$D_2$};
  \fill[DemandRed] (D4) circle (.024) node[above=3pt,figlabel] {$D_4$};
  \fill[CenterBlue] (E4) circle (.024) node[below=3pt,figlabel] {C trace end};
  \draw[GapPurple,line width=1.1pt] (E2)--(D2);
  \draw[black!35,line width=1.1pt] (D4)--(E4);
  \node[figlabel,align=center] at (.52,-.42)
    {one radial witness lies beyond the corresponding C trace};
\end{tikzpicture}%
}
\end{minipage}\hfill
\begin{minipage}[t]{.28\linewidth}
\centering
\resizebox{\linewidth}{!}{%
\begin{tikzpicture}[x=4.5cm,y=1cm,>=Latex]
  \node[panellabel,anchor=west] at (-.04,2.42) {(c) the case split};
  \node[figlabel,anchor=west] at (0,1.84) {first radial alternative};
  \draw[line width=.65pt] (0,1.38)--(1,1.38);
  \fill (0,1.38) circle (.017) node[left=1pt,figlabel] {$O$};
  \fill[CenterBlue] (.22,1.38) circle (.023) node[above=2pt,figlabel] {C trace};
  \fill[DemandRed] (.43,1.38) circle (.023) node[below=2pt,figlabel] {$D_2$};
  \draw[GapPurple,line width=1pt] (.22,1.38)--(.43,1.38);
  \node[figlabel,anchor=west] at (0,.73) {second radial alternative};
  \draw[line width=.65pt] (0,.27)--(1,.27);
  \fill (0,.27) circle (.017) node[left=1pt,figlabel] {$O$};
  \fill[CenterBlue] (.22,.27) circle (.023) node[above=2pt,figlabel] {C trace};
  \fill[DemandRed] (.43,.27) circle (.023) node[below=2pt,figlabel] {$D_4$};
  \draw[GapPurple,line width=1pt] (.22,.27)--(.43,.27);
  \node[figlabel,align=center] at (.5,-.36)
    {the red point lies beyond\\the corresponding center exit};
\end{tikzpicture}%
}
\end{minipage}
\caption{The CE2 short-ray terminal.  The two boundary gaps determine the
common pair $(p,q)$ and the center-forced radial witnesses $D_2,D_4$.  The
theorem proves that at least one of these witnesses lies beyond the
corresponding C-triangle trace and is uncovered.  The displayed ordering is
schematic.}
\label{fig:ce2-short-ray-certificate}
\end{figure}
